\section{From the certificate to pure absolute continuity}
\label{sec:spectral}

Fix an integer $d\ge3$ and a disorder parameter $\lambda>0$.
We prove a conditional transfer from a probability-$1/2$
finite-volume certificate to pure absolute continuity. The final
application checks the certificate at the couplings of
Theorem~\ref{thm:main}, using Proposition~\ref{prop:certificate} for
$d=3$ and Definition~\ref{def:threshold} together with
Proposition~\ref{prop:threshold} for $d\ge4$.

The finite-dimensional comparison bounds a bounded function of the
true scalar resolvent. A narrow spectral projection of the true
operator isolates the very near eigenvalues, whose effect is bounded
by a term proportional to its rank. The bound is valid for every
qualifying adaptive damping set. We then
pass to the infinite-volume boundary values and propagate positivity
to all sites at almost every fixed energy. The final argument
conditions on the operator with one site removed and excludes singular
mass, including mass on the exceptional energy set.
The operators considered below are bounded and self-adjoint, with
$\|H_{\lambda,\omega}\|\le2d+\lambda$. We use their spectral
resolution and functional calculus in the form of
\cite[Theorem~3.7]{Teschl2009}.

\subsection{Conditional spectral averaging}

We use own-site averaging twice: first to bound the expected number of
eigenvalues in a narrow true spectral window, and later to remove
spectral mass on an energy-null set fixed after freezing the other
potentials. We record its conditioning explicitly and prove the
precise form used below. For the general
rank-one theory and spectral-averaging background, see
\cite[Theorem~5 and Section~7]{SimonWolff1986} and
\cite[Theorem~1.1 and Proposition~3.1]{Marx2011}.

\begin{lemma}[Own-site averaging]\label{lem:spectral-averaging}
Let $\mathcal K$ be a Hilbert space, let $H^c$ be bounded and
self-adjoint on $\mathcal K$, and let $b\in\mathcal K$. On
$\CC\oplus\mathcal K$ set
\[
 H_s=\begin{pmatrix}s&b^*\\b&H^c\end{pmatrix},\qquad s\in\RR.
\]
Let $\mu_s$ be the spectral measure of $(1,0)$ for $H_s$.
If $s$ has a probability density $\rho\in L^\infty(\RR)$, then
for every Borel set $B\subset\RR$,
\begin{equation}\label{eq:spectral-averaging-bound}
 \int_{\RR}\rho(s)\mu_s(B)\,ds
 \le \|\rho\|_\infty |B|.
\end{equation}
This estimate holds conditionally when $H^c$ and $b$ are random and
independent of $s$.
\end{lemma}

\begin{proof}
For $\Im z>0$ define
\[
 M(z)=z+\langle b,(H^c-z)^{-1}b\rangle.
\]
Then $\Im M(z)>0$, and the Schur identity gives
\begin{equation}\label{eq:spectral-cavity}
 g_s(z):=\langle(1,0),(H_s-z)^{-1}(1,0)\rangle
       =\frac1{s-M(z)}.
\end{equation}
For any $a\in\RR$ and $v>0$,
\[
 \int_{\RR}\Im\frac1{s-a-iv}\,ds
 =\int_{\RR}\frac{v}{(s-a)^2+v^2}\,ds=\pi.
\]
Hence $\int\rho(s)\Im g_s(E+i\eta)\,ds
\le\pi\|\rho\|_\infty$. Multiply by a nonnegative continuous
compactly supported function $\varphi(E)$, integrate in $E$, and
divide by $\pi$. The spectral theorem and Fubini identify the
inner energy integral as the integral of the Poisson convolution of
$\varphi$ against $\mu_s$. That convolution converges uniformly to
$\varphi$ as $\eta\downarrow0$. Since $\mu_s(\RR)=1$, dominated
convergence gives
\[
 \int\rho(s)\int\varphi\,d\mu_s\,ds
 \le\|\rho\|_\infty\int\varphi(E)\,dE.
\]
Regularity of the two Borel measures proves
\eqref{eq:spectral-averaging-bound}. For a fixed value of the
conditioning variables the same calculation applies verbatim.
No cyclicity assumption on $(1,0)$ is needed.
\end{proof}

\subsection{The spectral readout}

\begin{proposition}[Spectral readout]\label{prop:certificate-readout}
Fix $d\ge3$, $\lambda>0$ and a nonempty bounded open interval $I$.
Let $H_m$ be the true Anderson operator on
$\mathbb T_m=(\ZZ/L_m\ZZ)^d$, $L_m=4^{20m}$, with independent
uniform diagonal values in $[-\lambda,\lambda]$ at every site,
and let $\eta_m\asymp4^{-m}$.
Suppose that for every fixed $E\in I$ and all sufficiently large
$m$, with probability at least $1/2$ there exist
$A\subset\mathbb T_m$ and numbers $r_i$, $i\in A$, such that
\[
 |A|\ge N_m/2,\qquad |r_i|\le B,\qquad
 0<p_-\le\Im r_i\le p_+,
\]
where $N_m=L_m^d$ and the displayed constants do not depend on
$m$ or $E$, and such that $(H_m-E-\eta_m\diag_A r)^{-1}$ exists
and its active block $\widehat G_A$, acting on $\ell^2(A)$, satisfies
\[
 \Im(\widehat G_A)_{ii}\ge q_*>0\qquad(i\in A),
\]
with $q_*$ independent of $m,E$. Every diagonal value of $H_m$, on
$A$ and on $A^c$, is already a true site value; here $A$ records the
sites retaining artificial damping. The certificate may depend on
all true variables and auxiliary randomness; no measurable
selection of $A,r$ is required. Then almost surely
\[
 I\subset\sigma(H_{\lambda,\omega}),\qquad
 \chi_I(H_{\lambda,\omega})\ell^2(\ZZ^d)
       \subset\mathcal H_{\mathrm{ac}}(H_{\lambda,\omega}).
\]
In particular, this spectral restriction is purely absolutely
continuous and $\langle e_0,\chi_I(H_{\lambda,\omega})e_0\rangle>0$.
\end{proposition}

\begin{proof}
Write $H=H_{\lambda,\omega}$. For $\Im z>0$, write
\[
 G(z)=(H-z)^{-1},\qquad
 G^{(m)}(z)=(H_m-z)^{-1},
\]
and put
\[
 q_x(z)=\Im G_{xx}(z),\qquad
 q_i^{(m)}(z)=\Im G_{ii}^{(m)}(z).
\]
When a diagonal resolvent entry has a finite complex boundary value,
write $q_x(E+i0)=\Im G_{xx}(E+i0)$. Our first objective is to prove
that, for Lebesgue-almost every fixed $E\in I$,
\begin{equation}\label{eq:spectral-all-positive}
 \PP\left(
 \begin{gathered}
 G_{xx}(E+i0)\text{ exists finitely and }q_x(E+i0)>0\\
 \text{for every }x\in\ZZ^d
 \end{gathered}
 \right)=1.
\end{equation}
The probability-one event here may depend on $E$. Once this assertion
is proved, the last part of the proof will use Fubini and own-site
averaging to obtain the almost-sure spectral statement on the whole
interval.

\paragraph{A true finite-volume resolvent bound.}
Fix $E\in I$ and suppress the subscripts on $\eta_m,N_m$.
We seek a lower bound for a bounded function of $q_i^{(m)}(E+i\eta)$
which is valid for the true operator even though the certificate is
chosen adaptively.

Let $\mathcal S_{m,E}$ be the event in the true potential variables
that a certificate exists. This event is Borel.
For each of the finitely many possible $A$, constrain the true
potential vector and $r$ to their compact ranges and impose
$|\det(H_m-E-\eta\diag_A r)|\ge1/j$, $j\in\NN$.
The inverse inequalities define a closed set there. Its projection
onto the potential coordinates is compact; taking the union over
$j$ and $A$ gives $\mathcal S_{m,E}$. An auxiliary construction
with success probability at least $1/2$ therefore implies
$\PP(\mathcal S_{m,E})\ge1/2$ for the true product law.

Almost surely, every true principal matrix of $H_m-E$ is invertible.
Indeed, each of its determinants is a nonzero polynomial in the
independent continuously distributed diagonal variables, and there
are finitely many subsets: the product of all diagonal variables
in a nonempty block has coefficient one in its determinant.
This event applies to an adaptively selected subset as well.
For any qualifying pair $(A,r)$, let $S_A$ be the real symmetric
Schur complement of $H_m-E$ after eliminating $A^c$. The active block of the
certified full inverse is
\[
 \widehat G_A=(S_A-\eta R)^{-1},\qquad
 R=\diag(r_i:i\in A).
\]
Define $G_A=(S_A-i\eta\Id)^{-1}$ on the same active space.
Here $\widehat G_A$ acts on $\ell^2(A)$, and dissipativity on that
space gives $\|\widehat G_A\|\le(\eta p_-)^{-1}$.
The inverse identity and the Ward identity give, respectively,
\[
 \widehat G_A-G_A=\widehat G_A\,\eta(R-i\Id)G_A,
 \qquad
 \Im\widehat G_A=\eta\widehat G_A^*\diag(\Im r)\widehat G_A.
\]
Set $C_R=1+(B+1)/p_-$. For $i\in A$, these identities imply
\[
 q_*\le\eta p_+\|\widehat G_Ae_i\|^2,\qquad
 \|\widehat G_Ae_i\|\le C_R\|G_Ae_i\|.
\]
Thus, with $c_*=q_*/(p_+C_R^2)>0$,
\begin{equation}\label{eq:spectral-subset-column}
 \eta\|G_Ae_i\|^2\ge c_*\qquad(i\in A).
\end{equation}

Put $T=(H_m-E)^{-1}$, and regard $B_A=P_ATP_A$ as an operator
on $\ell^2(A)$, where $P_A$ is coordinate projection. Block
inversion gives $S_A^{-1}=B_A$, and hence
\[
 G_A=P_A(H_m-E-i\eta P_A)^{-1}P_A=f_\eta(B_A),
 \qquad f_\eta(u)=\frac{u}{1-i\eta u}.
\]
Thus $G_A$ is an active-space intermediary: it is the compression of
an inverse damped on $A$. To compare it with the full scalar resolvent,
we separate the true spectral modes very close to $E$ before forming
a second compression. Those modes may make $T$ large, but the
difference they create after compression has rank at most their number.
Away from them, the
inverse-square weight is bounded by a multiple of the scalar
resolvent's Poisson weight.

For $0<\epsilon\le1$, define
\[
 Q=\chi_{[E-\epsilon\eta,E+\epsilon\eta]}(H_m),\quad
 r_Q=\operatorname{rank}Q,\quad
 T_c=T(\Id-Q),\quad B_c=P_AT_cP_A,\quad G_c=f_\eta(B_c).
\]
The identity
\[
 B_A-B_c=P_ATQ P_A
\]
shows that $\operatorname{rank}(B_A-B_c)\le r_Q$, for every set $A$.
For self-adjoint matrices $B,C$,
\[
 f_\eta(B)-f_\eta(C)
 =(\Id-i\eta B)^{-1}(B-C)(\Id-i\eta C)^{-1}.
\]
Thus $D=G_A-G_c$ has
\begin{equation}\label{eq:spectral-cut-rank}
 \operatorname{rank}D\le r_Q,\qquad
 \|D\|\le2/\eta,\qquad
 \|D\|_{\mathrm{HS}}^2\le4r_Q/\eta^2.
\end{equation}
Put $M_\epsilon=1+\epsilon^{-2}$.
Since $Q$ commutes with $T$, the matrix $T_c$ is self-adjoint,
and so is its compression $B_c$ on $\ell^2(A)$. Its own
functional calculus gives
\[
 G_c^*G_c=B_c^2(\Id+\eta^2B_c^2)^{-1}\le B_c^2.
\]
Together with contraction of the coordinate projection this implies
\begin{align}
 \eta\|G_ce_i\|^2
 &\le\eta\|B_ce_i\|^2\le\eta\|T_ce_i\|^2\notag\\
 &\le M_\epsilon q_i^{(m)}(E+i\eta),\qquad i\in A.
 \label{eq:spectral-cut-column}
\end{align}
For the last inequality, outside the cut interval one has
$\eta/(u-E)^2\le M_\epsilon\eta/((u-E)^2+\eta^2)$.
The argument neither commutes $P_A$ past $Q$ nor uses an
operator-monotonicity assertion about $f_\eta$.

Define the bounded continuous function
$\Phi(u)=\min\{2M_\epsilon u,c_*\}$ for $u\ge0$.
Equations~\eqref{eq:spectral-subset-column} and
\eqref{eq:spectral-cut-column}, together with the squared triangle
inequality, give
\[
 \Phi(q_i^{(m)})\ge c_*-2\eta\|De_i\|^2\qquad(i\in A).
\]
Summing and using \eqref{eq:spectral-cut-rank} yields
\[
 \frac1N\sum_i\Phi(q_i^{(m)})
 \ge c_*/2-\frac{8r_Q}{\eta N}.
\]
Every qualifying certificate implies this same inequality for the
true scalar resolvent and the true rank $r_Q$. Thus it holds on
the existential event $\mathcal S_{m,E}$ without choosing a
certificate measurably. Outside that event the left side is
nonnegative, so pointwise it is at least
$\frac{c_*}{2}\mathbf1_{\mathcal S_{m,E}}-8r_Q/(\eta N)$.
Taking expectations under the unconditioned true product law gives
\begin{equation}\label{eq:spectral-truncated-average}
 \EE\frac1N\sum_i\Phi(q_i^{(m)})
 \ge\frac{c_*}{2}\PP(\mathcal S_{m,E})
       -\frac{8\EE r_Q}{\eta N}.
\end{equation}

Under this law, at each site the true diagonal value is uniform on
$[-\lambda,\lambda]$ and independent of the other site values.
Conditioning on those other values, applying
Lemma~\ref{lem:spectral-averaging}, and summing gives the finite-volume
Wegner estimate~\cite{Wegner1981}
\[
 \EE r_Q\le\frac{N}{2\lambda}(2\epsilon\eta)
         =\frac{\epsilon\eta N}{\lambda}.
\]
The rank estimate is under the unconditioned true law; the success
event enters only through its probability in
\eqref{eq:spectral-truncated-average}. Choose
$0<\epsilon\le\min\{1,c_*\lambda/64\}$, fixed as $m$ varies.
Using $\PP(\mathcal S_{m,E})\ge1/2$, for all sufficiently large
$m$ Equation~\eqref{eq:spectral-truncated-average}
therefore gives
\[
 \EE\frac1N\sum_i\Phi(q_i^{(m)}(E+i\eta_m))\ge c_*/8.
\]
The expression being averaged now depends only on the true
potentials. Translation invariance of that true product law
on the torus yields
\begin{equation}\label{eq:spectral-finite-root}
 \EE\Phi(q_0^{(m)}(E+i\eta_m))\ge c_*/8.
\end{equation}
The construction of $A$ need not be translation equivariant.

\paragraph{Infinite volume and the first boundary conclusion.}
We next transfer \eqref{eq:spectral-finite-root} to the infinite
lattice. Couple the potentials on a root cube of radius comparable
to $L_m$, less than half the torus side, to the infinite iid model.
Cutting the boundary bonds of this cube gives the same root
component in the two models. Exponential conjugation of a
scalar-damped nearest-neighbor resolvent gives
\[
 |G(x,y;E+i\eta)|\le C\eta^{-1}e^{-c\eta\dist(x,y)}.
\]
Indeed, conjugation by $e^{a\eta\dist(x,\cdot)}$ changes the
hopping in operator norm by at most $Ca\eta$; a Neumann series
and $\|G\|\le\eta^{-1}$ prove the estimate for sufficiently small
fixed $a$. Truncated weights give the same conclusion on the
infinite lattice. The bound also holds after boundary bonds are
removed. A cube of side comparable to $L_m$ has $2d$ faces,
each containing $O_d(L_m^{d-1})$ crossing bonds. Resolvent
identity across these bonds, with the root a distance
$\asymp L_m$ from every face, gives
\begin{equation}\label{eq:spectral-volume-comparison}
 |G_{00}^{(m)}(E+i\eta_m)-G_{00}(E+i\eta_m)|
 \le C_d L_m^{d-1}\eta_m^{-2}e^{-c_d\eta_mL_m}=o(1).
\end{equation}
Indeed, $\eta_mL_m\asymp4^{19m}$, whereas the logarithm of
the prefactor is $O_d(m)$.
This estimate is uniform in the potential. Since $\Phi$ is bounded
and Lipschitz, it follows that
\begin{equation}\label{eq:spectral-infinite-root}
 \liminf_m\EE\Phi(q_0(E+i\eta_m))\ge c_*/8.
\end{equation}

Each diagonal resolvent entry is the Borel transform of a finite measure,
and thus has a finite complex boundary value for Lebesgue-almost
every energy; see \cite[Corollary~3.25]{Teschl2009}.
Fubini gives a deterministic full-measure subset of $I$ on which
these limits exist almost surely. Bounded convergence in
\eqref{eq:spectral-infinite-root} shows that, for almost every
fixed $E\in I$,
\begin{equation}\label{eq:spectral-positive-probability}
 \EE\Phi(q_0(E+i0))\ge c_*/8,
 \qquad \PP(q_0(E+i0)>0)>0.
\end{equation}
All subsequent energy limits concern the true resolvent. In
particular, no joint choice of the auxiliary construction over
uncountably many energies is required.

\paragraph{Boundary nonvanishing and propagation of positivity.}
The shortest-path observation and rank-one comparison below have
precedents in Jak\v{s}i\'c and Last
\cite[author-preprint Section~1.1 and proof of Theorem~2.4]{JaksicLast2000}.
We give the boundary-value argument used here. We first claim that,
for every fixed realization and pair $x\ne y$, $G_{xy}(E+i0)$ exists
finitely and is nonzero for Lebesgue-almost every $E$. This form will
let us apply Fubini before returning to a fixed energy.

To prove the claim, real symmetry gives
\[
 G_{xy}(z)=\tfrac14\bigl(
 \langle e_x+e_y,G(z)(e_x+e_y)\rangle
 -\langle e_x-e_y,G(z)(e_x-e_y)\rangle\bigr).
\]
Both quadratic forms are Herglotz functions, so $G_{xy}$ has finite
boundary values a.e. To prove nonvanishing, let $\ell=|x-y|_1$.
The first nonzero coefficient of
the convergent expansion
\[
 (H-z)^{-1}=-\sum_{n\ge0}H^n z^{-n-1}
\]
in the $(x,y)$ entry occurs at $n=\ell$. A product of fewer
than $\ell$ nearest-neighbor steps and diagonal factors cannot
join $x$ to $y$; at length $\ell$ every factor must be a hopping
step on a shortest path. Thus $(H^\ell)_{xy}$ equals
$(-1)^\ell$ times the positive number of shortest lattice paths
and $G_{xy}$ is not identically zero in any dimension considered.

Write each quadratic Herglotz function $m_\pm$ as
$m_\pm=i(1+b_\pm)/(1-b_\pm)$ with $b_\pm$ bounded analytic
and $|b_\pm|<1$. Their difference is a quotient of bounded
analytic functions with nonzero numerator. A nonzero bounded
analytic function $f$ on the upper half-plane cannot have zero
vertical boundary values on a positive-measure set. For
completeness, normalize $|f|\le1$ and choose $z_0$ with
$f(z_0)\ne0$. Jensen's inequality after mapping
$\{\Im z>\eta\}$ to the disk gives
\[
 \int_{\RR}\frac{\Im z_0-\eta}
 {\pi((x-\Re z_0)^2+(\Im z_0-\eta)^2)}
       [-\log|f(x+i\eta)|]\,dx
 \le-\log|f(z_0)|
\]
for $0<\eta<\Im z_0$. To obtain this boundary inequality from
Jensen on smaller circles, use Fatou's lemma and the continuity
of $f$ across the line $\Im z=\eta$; the single boundary point
corresponding to infinity has zero harmonic measure.
Another application of Fatou's lemma as $\eta\downarrow0$
rules out a positive-measure zero-boundary set. If the finite
boundary value of $G_{xy}$ were zero, the bounded-analytic
numerator would tend to zero too, since the denominator factors
$1-b_\pm$ are bounded. Applying the uniqueness statement to
that numerator proves the claimed nonvanishing.

For each realization, intersect the boundary statements over the
countable set of sites and pairs, and then apply Fubini. Together
with \eqref{eq:spectral-positive-probability}, this gives a deterministic
full-measure set $I_0\subset I$ such that for every fixed $E\in I_0$,
\eqref{eq:spectral-positive-probability} holds and, almost surely, all
diagonal and off-diagonal boundary values are finite and every
off-diagonal value for distinct sites is nonzero. Fix such an $E$
for the rest of this propagation argument. The event that some site
has $q_x(E+i0)>0$ is translation invariant and has positive
probability. Ergodicity of the iid translation action makes its
probability one. One elementary proof of the required ergodicity
approximates an invariant event by a finite-coordinate cylinder
event, translates the latter to make the two cylinder events
independent, and lets the approximation error tend to zero.

Fix an ordered pair $x\ne y$ and condition on all potentials
except the true diagonal value $s_x$. For almost every such
complement, the finite-boundary and nonzero-off-diagonal
properties hold for almost every $s_x\in[-\lambda,\lambda]$.
Call these own-site values admissible. This full-measure set
may depend on the complement and on $E$; the argument requires
admissibility only for almost every own-site value at this fixed energy.
There cannot be two distinct admissible values $s_1,s_2$ for
which $q_x>0$ and $q_y=0$. To see this, evaluate at $s_1$ and set
\[
 g=G_{xx}(E+i0),\quad a=G_{xy}(E+i0),\quad h=s_2-s_1.
\]
The boundary imaginary-part matrix on $\{x,y\}$ is positive
semidefinite, being a limit of such matrices.
Its $y$-diagonal is zero, so its $(x,y)$ entry vanishes. Hence
$a$ is real; it is nonzero by the preceding argument. Moreover,
$\Im g>0$, so $1+hg\ne0$ for every real $h$.
The upper-half-plane rank-one identity therefore passes to the
boundary and gives
\[
 G_{yy}^{(s_2)}=G_{yy}^{(s_1)}-\frac{h a^2}{1+hg},
 \qquad
 q_y^{(s_2)}=\frac{h^2a^2\Im g}{|1+hg|^2}>0,
\]
a contradiction. The set of admissible $s_x$ giving
$q_x>0,q_y=0$ therefore contains at most one point and has zero
conditional probability under the uniform own-site law. The
inadmissible values already form a conditional null set. Thus
\[
 \PP(q_x(E+i0)>0,\ q_y(E+i0)=0)=0.
\]
Countably many pairs, combined with almost-sure existence of a
positive site, prove \eqref{eq:spectral-all-positive} for every
$E\in I_0$. This completes the fixed-energy part of the proof.

\paragraph{Exclusion of singular mass.}
We now use the fixed-energy assertion \eqref{eq:spectral-all-positive}
to prove an almost-sure statement about the spectral measure on $I$.
Positivity a.e. in energy is not by itself sufficient: a spectral
measure could also have singular mass on an exceptional null set.
After fixing the other site potentials, we construct an energy set on
which every value of the remaining site potential has a finite diagonal
boundary value. Its complement can then be used in conditional spectral
averaging.
This conditional spectral-averaging step is the purity principle of
Aizenman and Warzel~\cite[Proposition~2.2]{AizenmanWarzel2013}, credited
there to Mira Shamis via Simon--Wolff's arguments; compare also
Jak\v{s}i\'c and Last~\cite[author-preprint Section~4]{JaksicLast2000}.

Fix a site $x$, and write $\overline\omega$ for the potentials at
the other sites. Let $H^c_{\overline\omega}$ be the compression with
site $x$ removed, and let $b$ be the hopping vector from $x$ to the
remaining sites. Both are independent of the true diagonal value
$s$ at $x$. For a real value $s$ at that site, the cavity identity
\eqref{eq:spectral-cavity} becomes
\[
 G_{xx}^{(s)}(z)=\frac1{s-M_{\overline\omega}(z)},\qquad
 M_{\overline\omega}(z)
 =z+\langle b,(H^c_{\overline\omega}-z)^{-1}b\rangle.
\]
The function $M_{\overline\omega}$ has finite boundary values
for almost every energy, since its second term is a Borel
transform of a finite positive measure. Its finite boundary values
have nonnegative imaginary part. By
\eqref{eq:spectral-all-positive} and Fubini, for almost every
$\overline\omega$ and almost every $E\in I$, the diagonal boundary
value exists finitely and has positive imaginary part for almost every
$s\in[-\lambda,\lambda]$.
If the finite value
$M_{\overline\omega}(E+i0)$ were real, the cavity formula would
give zero imaginary boundary part for every
$s\ne M_{\overline\omega}(E+i0)$, a contradiction. Thus the set
\[
 S_{\overline\omega}
 =\{E\in I:M_{\overline\omega}(E+i0)\text{ exists finitely and }
                      \Im M_{\overline\omega}(E+i0)>0\}
\]
has full Lebesgue measure in $I$ for almost every complement.
The boundary condition defining $S_{\overline\omega}$ is Borel and
jointly measurable in $(\overline\omega,E)$: at positive height the
resolvents are jointly measurable and continuous in height, so boundary
existence and its bounds can be expressed using rational heights.
The set depends only on the complementary configuration, and hence
is independent of $s$. At each point of $S_{\overline\omega}$, the
positive imaginary part of $M_{\overline\omega}(E+i0)$ keeps the cavity
denominator nonzero for every real $s$. Thus the diagonal resolvent has
a finite boundary value there for every $s$.

We spell out the measure-theoretic consequence. Let $\mu$ be a
finite positive measure with Borel transform $g$, and let $S$ be
a Borel set on which
$\limsup_{\eta\downarrow0}\Im g(E+i\eta)<\infty$. Put
\[
 S_k=\{E\in S:\sup_{0<\eta<1/k}\Im g(E+i\eta)\le k\}.
\]
These sets are Borel, since continuity at positive height allows the
supremum to be taken over rational heights. Then $S=\bigcup_k S_k$.
For $E\in S_k$ and $0<r<1/k$,
\begin{equation}\label{eq:spectral-poisson-cover}
 \mu((E-r,E+r))\le2r\Im g(E+ir)\le2kr.
\end{equation}
Let $B\subset S_k$ be a Lebesgue-null Borel set. For each
$\varepsilon>0$, cover $B$ by open intervals $J_j$ with
$\sum_j|J_j|<\varepsilon$ and $|J_j|<1/k$.
Discard intervals missing $B$ and choose $E_j\in B\cap J_j$.
The intervals $(E_j-|J_j|,E_j+|J_j|)$ still cover $B$, so
Equation~\eqref{eq:spectral-poisson-cover} gives
\[
 \mu(B)\le\sum_j\mu((E_j-|J_j|,E_j+|J_j|))
 \le2k\sum_j|J_j|<2k\varepsilon.
\]
Letting $\varepsilon\downarrow0$ proves $\mu|_{S_k}\ll dE$, and hence
$\mu|_S\ll dE$. This proves directly the bounded-boundary
criterion in the Borel-set form we need; compare
\cite[Theorem~3.23 and Corollary~3.24]{Teschl2009}.

Apply this deterministic argument, for each $s$, to its site measure $\mu_s$
on $S_{\overline\omega}$. The bounded-Poisson subsets $S_k$
may depend on $s$, but the larger set $S_{\overline\omega}$
does not. Therefore
\[
 \mu_s|_{S_{\overline\omega}}\ll dE\quad\text{for every }s.
\]
For its complementary Borel null set
$B_{\overline\omega}=I\setminus S_{\overline\omega}$,
conditional spectral averaging gives
\[
 \int_{-\lambda}^{\lambda}\frac1{2\lambda}
       \mu_s(B_{\overline\omega})\,ds
 \le\frac1{2\lambda}|B_{\overline\omega}|=0.
\]
For almost every complement, this proves
$\mu_s(B_{\overline\omega})=0$ for almost every own-site value
$s\in[-\lambda,\lambda]$.
Together with absolute continuity on $S_{\overline\omega}$ for every
$s$, it follows that $\mu_x|_I\ll dE$ almost surely under the true
product law. Countably intersecting over $x$ gives simultaneous
absolute continuity of all basis measures on $I$.

The vectors $\chi_I(H)e_x$ consequently belong to
$\mathcal H_{\mathrm{ac}}(H)$. Their linear span is dense in
$\chi_I(H)\ell^2(\ZZ^d)$, and the absolutely continuous
subspace is closed. This proves the spectral-subspace inclusion.
Finally, \eqref{eq:spectral-all-positive} and Fubini give
$q_0(E+i0)>0$ for a.e. $E\in I$ on a probability-one event.
If some point of $I$ lay outside the closed spectrum, it would
have an open neighborhood in $I$ disjoint from the spectrum.
The resolvent is real there, so its diagonal imaginary boundary
parts vanish throughout that neighborhood, a contradiction.
Thus $I\subset\sigma(H)$ as well.
Moreover $\mu_0(I)>0$: if $\mu_0(I)=0$, the Poisson integral
of $\mu_0$ tends to zero at each interior point of $I$, since
that point has positive distance from $\RR\setminus I$.
This contradicts $q_0(E+i0)>0$ for almost every $E\in I$.
In particular $\chi_I(H)\ne0$.
\end{proof}

\begin{proof}[Proof of Theorem~\ref{thm:main}]
For $d=3$, fix the interval $I_3=(-6+e_2/2,-6+e_2)$ and set
$\lambda_3=\lambda_0(3)>0$ from Proposition~\ref{prop:certificate}.
For each fixed $0<\lambda<\lambda_3$, that proposition supplies the
certificate with probability tending to one, hence at least $1/2$
for all sufficiently large $m$, and with constants uniform in
$E\in I_3$. Proposition~\ref{prop:certificate-readout} therefore proves
the two spectral assertions on $I_3$.

For $d\ge4$, set
\[
 I_d=-2d+(1/200,1/100),\qquad
 \lambda_d=\sqrt3\,2^{-n_*(d)}>0,
\]
where $n_*(d)$ is specified in Definition~\ref{def:threshold}.
For each fixed $0<\lambda<\lambda_d$, put
\[
 n_0=\lfloor\log_4(3/\lambda^2)\rfloor\ge n_*(d),
 \qquad u=(\lambda^2/3)4^{n_0}\in(1/4,1].
\]
The existence assertion in Proposition~\ref{prop:threshold} follows
from the construction proving Proposition~\ref{prop:certificate}.
Its final assertion applies the defining feasibility criterion
directly: for this coupling and each $m\ge4n_0$ and $E\in I_d$,
the true-variable certificate event has probability at least $1/2$.
With $\eta_m=s_m=u4^{-m}$, its fixed bounds
$B=p_+=n_*(d)$ and $p_-=q_*=1/n_*(d)$ satisfy precisely the
hypotheses of Proposition~\ref{prop:certificate-readout} on $I_d$.
Hence almost surely $I_d$ is contained in the spectrum, and its
spectral restriction is purely absolutely continuous and has nonzero
weight. In either dimension case, conjugating by
$\psi(x)\mapsto(-1)^{x_1+\cdots+x_d}\psi(x)$ gives
the identical conclusion for positive hopping and leaves every
iid diagonal variable unchanged.
The interval is independent of $\lambda$, and the probability-one
event is asserted separately for each fixed disorder parameter.
This proves the stated spectral conclusion; no diffusion
assertion is used or deduced.
\end{proof}
